%==============================================================================
% Chapter 2 — Engineering Foundations & Infrastructure
% Curriculum: Phase 1, Chapter 2 of docs/AI_ML_Master_Checklist.md
% Companion code: projects/02_engineering_foundations/
%
% Section skeleton only. Figures for this chapter go in theory/figures/.
%==============================================================================
\chapter{Engineering Foundations \& Infrastructure}
\label{ch:engineering-foundations}

Chapter~\ref{ch:math-foundations} covered the language the models are written
in.\projectref{projects/02\_engineering\_foundations/} This chapter covers the machinery that lets them be built twice and get the
same answer. It sits in Phase~1 deliberately: dependency management, data
versioning, branch protection and continuous integration are cheap to establish
now and expensive to retrofit once eight chapters of code already exist. Every
chapter from~\ref{ch:classical-ml} onward assumes the standards fixed here.

The chapter divides into a theory half --- the design principles and the
reasoning behind them --- and an applied half, which is the actual toolchain
configured in this repository.

%==============================================================================
\section{Systems Design Principles}
\label{sec:systems-design}

\subsection{Separation of Concerns}

\subsection{Dependency Direction}

\subsection{Idempotency}

\subsection{The Reproducibility Contract}

%==============================================================================
\section{Monorepo Architecture}
\label{sec:monorepo}

\subsection{Why Theory, Library Code and Projects Share a Repository}

\subsection{Shared-Library versus Per-Project Dependency Boundaries}

\subsection{The Cost Model: One Repository versus Many}

%==============================================================================
\section{Isolated Environments}
\label{sec:isolated-environments}

\subsection{Virtual Environments and the Interpreter Boundary}

\subsection{Lockfiles versus Loose Version Ranges}

\subsection{Why a Global \texttt{requirements.txt} Fails}

%==============================================================================
\section{Dependency Management with \texttt{uv}}
\label{sec:uv}

\subsection{Project Initialisation}

\subsection{The Core Workflow: \texttt{sync}, \texttt{add}, \texttt{run}}

\subsection{Resolving and Committing \texttt{uv.lock}}

\subsection{Pinning the Interpreter with \texttt{.python-version}}

\subsection{Per-Project \texttt{pyproject.toml} for Isolation}

%==============================================================================
\section{Data Version Control}
\label{sec:dvc}

\subsection{Initialising DVC}

\subsection{Tracking Datasets as Lightweight Pointers}

\subsection{Configuring a Remote}

\subsection{The \texttt{push}/\texttt{pull} Round Trip}

%==============================================================================
\section{Git Workflow and Branch Protection}
\label{sec:branch-protection}

\subsection{Trunk-Based Development and Feature Branches}

\subsection{Required Reviews and Status Checks}

\subsection{Linear History}

%==============================================================================
\section{Continuous Integration and Delivery}
\label{sec:cicd}

\subsection{Workflow Anatomy: Triggers, Jobs, Steps}

\subsection{Path Filters}

\subsection{Automated Compilation of This Book}

\subsection{Testing and Linting Gates}

\subsection{Artifact Upload}

%==============================================================================
\section{Core Software Engineering Practice}
\label{sec:software-engineering}

\subsection{Unit Testing with \texttt{pytest}}

\subsection{Fixtures and Parametrisation}

\subsection{Type Hints and Docstrings as a Baseline}
